% =====================================================================
% II. PLANTEAMIENTO DEL PROBLEMA
% =====================================================================
\section{Planteamiento del Problema}

En la pr\'actica, la digitalizaci\'on de partituras contin\'ua siendo un
proceso predominantemente reactivo y dependiente del esfuerzo humano.
Aunque el reconocimiento \'optico de m\'usica ha evolucionado desde sistemas
basados en reglas hasta \emph{pipelines} de aprendizaje profundo capaces de
transcribir p\'aginas completas de forma autom\'atica
\cite{omr2020}\cite{calvo2018}\cite{riosvila2026}, su salida a\'un comete
errores frecuentes que impiden usarla directamente \cite{calvo2018}: alturas
o duraciones incorrectas, compases desbalanceados, armaduras inconsistentes
y s\'imbolos omitidos. Por eso, quien digitaliza termina corrigiendo
s\'imbolo por s\'imbolo en un editor convencional, con un costo de tiempo
cercano al de la transcripci\'on manual, o bien acepta resultados sin
verificaci\'on alguna. En el caso de los manuscritos modernos esta situaci\'on se
agrava: la irregularidad del trazo, la inclinaci\'on de las plicas y la
variabilidad geom\'etrica de los s\'imbolos degradan a\'un m\'as la
confiabilidad del reconocimiento autom\'atico \cite{ricordi2024}.

Esta condici\'on evidencia una brecha disciplinar: los modelos de
reconocimiento se estudian predominantemente desde la precisi\'on de su
salida autom\'atica, mientras que su integraci\'on en el trabajo pr\'actico
de un transcriptor ha recibido mucha menos atenci\'on \cite{rizo2023}. El
objeto de estudio de este proyecto es, por tanto, el dise\~no de un sistema
de informaci\'on que resuelva la fragmentaci\'on actual entre tres
dimensiones que hoy operan de manera aislada: (i) los \emph{modelos de
reconocimiento OMR}, evaluados casi exclusivamente mediante m\'etricas de
precisi\'on simb\'olica; (ii) las t\'ecnicas de \emph{validaci\'on musical},
capaces de detectar errores estructurales de m\'etrica, duraciones y
armaduras mediante reglas de teor\'ia musical sin intervenci\'on humana
\cite{cuthbert2010}; y (iii) los \emph{mecanismos de interacci\'on humana},
que permiten al usuario corregir la transcripci\'on sobre la imagen
original, pero cuyas correcciones, en los sistemas actuales, se descartan
sin alimentar con ellas la mejora progresiva del sistema mediante
aprendizaje activo, enfoque pr\'acticamente inexplorado en OMR y con
resultados iniciales incluso contradictorios \cite{amershi2014}%
\cite{settles2009}\cite{mlsp2025}.

A nivel de ingenier\'ia, el problema se concreta en la inexistencia de una
plataforma tecnol\'ogica integrada que articule estas dimensiones. En el
marco experimental de este estudio, las variables se definen de la
siguiente manera: la \emph{Variable Independiente} es la modalidad del
flujo de digitalizaci\'on aplicada a cada entrada, entendida como la
comparaci\'on entre una l\'inea base OMR no asistida y el flujo asistido
con validaci\'on musical y aprendizaje activo, sobre partituras impresas
estructuradas y manuscritas modernas; la \emph{Variable Mediadora} es el
procesamiento anal\'itico de la salida, es decir, la detecci\'on de
anomal\'ias mediante reglas de teor\'ia musical y la selecci\'on activa de
muestras para el reentrenamiento; y las \emph{Variables Dependientes} son
el esfuerzo de correcci\'on humana, cuantificado en tiempo de edici\'on e
intervenciones manuales por comp\'as, y la fidelidad de la transcripci\'on
final (MusicXML/MIDI), medida mediante la tasa de error simb\'olico
(\emph{Symbol Error Rate}, SER).

En el contexto regional, esta necesidad afecta a archivos, academias y
agrupaciones musicales que conservan colecciones de partituras impresas y
escritas a mano cuya transcripci\'on manual representa costos de tiempo
considerables, y que cuentan \'unicamente con editores convencionales, sin
mecanismos de verificaci\'on autom\'atica ni de mejora continua. La
experiencia reciente de digitalizaci\'on de colecciones musicales,
abordada con metodolog\'ias de anotaci\'on y clasificaci\'on autom\'atica
\cite{ricordi2024}, confirma que el problema es real y abordable desde la
ingenier\'ia. A ello se suma que los conjuntos de datos p\'ublicos con
\emph{ground truth}, como PrIMuS/Camera-PrIMuS \cite{calvo2018}, Sheet
Music Benchmark \cite{smb2026} y MUSCIMA++ \cite{muscima2017}, permiten
evaluar la soluci\'on propuesta de forma cuantitativa y reproducible,
siguiendo los protocolos de la literatura de OMR.

\textbf{Pregunta de Investigaci\'on:} \textquestiondown De qu\'e manera el
dise\~no e implementaci\'on de una plataforma de digitalizaci\'on asistida
de partituras, que integre reconocimiento \'optico de m\'usica, validaci\'on
autom\'atica mediante reglas de teor\'ia musical y un ciclo de aprendizaje
activo a partir de las correcciones del usuario, permite reducir el
esfuerzo de revisi\'on humana y mejorar la fidelidad de la transcripci\'on
digital (MusicXML/MIDI), medida en tasa de error simb\'olico y tiempo de
edici\'on, en partituras impresas y manuscritos modernos de complejidad
monof\'onica o piano simple?
